\begin{lemma}
\label{editorial-source-lemma-support-quotient}
Let $R$ be a ring, let $M$ be an $R$-module, and let $I\subseteq R$ be an ideal.
\begin{enumerate}
\item If $M$ is finite then the support
of $M/IM$ is $\text{Supp}(M) \cap V(I)$.
\item If $N \subset M$, then $\text{Supp}(N) \subset
\text{Supp}(M)$.
\item If $Q$ is a quotient module of $M$ then $\text{Supp}(Q) \subset
\text{Supp}(M)$.
\item If $0 \to N \to M \to Q \to 0$ is a short exact sequence
then $\text{Supp}(M) = \text{Supp}(Q) \cup \text{Supp}(N)$.
\end{enumerate}
\end{lemma}

\begin{proof}
Localization preserves the original injections, surjections
and short exact sequences. For (2), the map
$N_{\mathfrak p}\to M_{\mathfrak p}$ is injective, so
the first module can be nonzero only if the second is.
For (3), the map $M_{\mathfrak p}\to Q_{\mathfrak p}$
is surjective, so the same conclusion holds for a nonzero
quotient. In (4), the exact localized sequence shows
$M_{\mathfrak p}=0$ if and only if both
$N_{\mathfrak p}=0$ and $Q_{\mathfrak p}=0$, which is
precisely the stated union formula.

For (1), retain the exact quotient comparison
$$
(M/IM)_{\mathfrak p}\longrightarrow M_{\mathfrak p}/IM_{\mathfrak p},
\qquad (m+IM)/s\longmapsto m/s+IM_{\mathfrak p}.
$$
Its inverse sends $m/s+IM_{\mathfrak p}$ to $(m+IM)/s$.
The original quotient map localizes to a surjection with
kernel $(IM)_{\mathfrak p}=IM_{\mathfrak p}$: the equality
follows by writing finite sums of products over a common
denominator. This proves well-definedness and both
composites of the displayed maps. If $M_{\mathfrak p}=0$,
the quotient is zero. If some $a\in I$ lies outside
$\mathfrak p$, it is a unit in $R_{\mathfrak p}$ and
$IM_{\mathfrak p}=M_{\mathfrak p}$, so the quotient is
again zero. Conversely, if $M_{\mathfrak p}\ne0$ and
$I\subseteq\mathfrak p$, then $IR_{\mathfrak p}$ lies
in the maximal ideal of the local ring. The module
$M_{\mathfrak p}$ is finite, so Nakayama's
Lemma \ref{algebra-lemma-NAK} prevents
$M_{\mathfrak p}/IM_{\mathfrak p}$ from being zero.
These cases prove the exact intersection formula.
\end{proof}